% !TeX root = ../../Book1.tex
% 纲要标题：### A.1 Zorn 引理与选择公理
% 纲要位置：计划纲要.md，第 813 行。
% BEGIN APPENDIX OUTLINE SYNC
% #### A.1.1 选择函数与选择公理
% #### A.1.2 偏序集与 Zorn 引理
% #### A.1.3 极大对象与基的扩充
% #### A.1.4 良序与超限归纳
% #### A.1.5 无限集合与有限支集计数
% END APPENDIX OUTLINE SYNC

\section{Zorn 引理与选择公理}
\label{b1:sec:A01}

在有限维向量空间中，从一个线性无关组出发，不断加入线性生成空间以外的向量，经过有限步便能得到一组基。若空间没有预先给定的有限维数，这个办法就留下了一个问题：即使每一步都能继续添加，怎样知道存在一个已经不能再扩大的线性无关集？Zorn 引理回答的正是这一类存在性问题。本节先说明它所依据的选择原则，再用基的扩充把论证中的每一步写清楚。

\ProseParagraphBreak
本节前三个小节讨论选择公理、Zorn 引理和基的扩充，足以支撑正文中极大理想及基的存在性证明。后两个小节的良序与无限基数计数则用于任意秩自由交换群、无限基和超越基的基数论证。

\subsection{选择函数与选择公理}
\label{b1:subsec:A01:01}

设有一族非空集合 \((X_i)_{i\in I}\)。所谓从每个集合中选一个元素，是要给出一个以 \(I\) 为定义域的函数：指标为 \(i\) 时，函数值必须属于 \(X_i\)。各个 \(X_i\) 可以相交，也可以相同；指标说明这一次选择对应于哪一个集合。

\begin{definition}[选择公理]\label{b1:def:A-choice}
对集合族 \((X_i)_{i\in I}\)，若函数
\[
c:I\longrightarrow\bigcup_{i\in I}X_i
\]
满足对每个 \(i\in I\) 都有 \(c(i)\in X_i\)，则称 \(c\) 为该族的\term[xuanze-hanshu]{选择函数}[choice function]。\term[xuanze-gongli]{选择公理}[axiom of choice]断言：每个由非空集合组成的、以集合为指标的集合族都有选择函数。
\end{definition}

例如 \(X_1=\{1,2\}\)、\(X_2=\{2,3\}\) 时，可以取 \(c(1)=1\)、\(c(2)=2\)。对有限指标集，总能用有限归纳证明选择函数存在：已有前几个集合的选择，再从下一个非空集合中取一个元素即可。这种有限次选取不需要另加选择公理。

\ProseParagraphBreak
无限族也可能有明确的选择方法。若每个 \(X_i\) 都是 \(\mathbb N\) 的非空子集，便可统一规定 \(c(i)=\min X_i\)，无论指标集有多大都没有困难。选择公理所增加的内容，是在集合族没有这种已知统一规则时，仍保证同时选择的函数存在。这里要区分两个量词次序：
\[
\forall i\in I\;\exists x\in X_i,
\qquad
\exists c\;\forall i\in I,\quad c(i)\in X_i.
\]
前者只说各个集合非空，后者要求一个在所有指标上同时有定义的函数。即使指标是自然数，把“从 \(X_0\) 选一个，再从 \(X_1\) 选一个，如此继续”写成一句话，也没有自动完成证明。递归构造需要事先给定每一步由哪些数据决定下一取值的规则；“从非空集合中任选一个”尚不是这种已定义的规则，而 \(c(i)=\min X_i\) 才是。对任意非空集合族同时保证选择函数存在，正是所采用公理的内容。

\ProseParagraphBreak
积集 \(\prod_{i\in I}X_i\) 的元素就是满足 \(x_i\in X_i\) 的全部坐标族 \((x_i)_{i\in I}\)。所以选择公理也可以写成：非空集合的积非空。指标集为空时，唯一的空函数给出空积的唯一元素。本书采用选择公理；它是一项存在性原则，并不提供计算各个 \(c(i)\) 的算法。

\ProseParagraphBreak
一个常见用法是为满射 \(q:X\rightarrow Y\) 选取集合意义下的截面。对每个 \(y\in Y\)，纤维 \(q^{-1}(y)\) 非空；在其中选一点 \(s(y)\)，便得到 \(s:Y\rightarrow X\) 和 \(qs=\operatorname{id}_Y\)。这只选出了元素，并没有保证选择保持额外结构；即使 \(q\) 是群或环的同态，所得集合截面也未必是同态。例如商同态 \(\mathbb Z\rightarrow\mathbb Z/2\mathbb Z\) 没有群同态截面：若存在这样的截面，\(\overline1\) 的像 \(z\) 应满足 \(2z=0\)，故 \(z=0\)，却又必须映回 \(\overline1\)，矛盾。

\subsection{偏序集与 Zorn 引理}
\label{b1:subsec:A01:04}

研究能否继续扩大的对象时，需要先规定“扩大”是什么意思。对子集采用包含关系；对部分定义的映射，则要求新映射保留旧定义域上的全部取值。不同对象可能互不包含，因此这里通常不能把所有对象排成一列大小次序。

\begin{definition}\label{b1:def:poset-chain}
设 \(\leq\) 是集合 \(P\) 上的二元关系，即 \(P\times P\) 的一个子集，并以 \(x\leq y\) 表示 \((x,y)\) 属于该关系。若它满足以下三条，则称 \(\leq\) 为\term[pianxu]{偏序}[partial order]，称 \((P,\leq)\) 为偏序集：
\begin{enumerate}
\item 自反性：每个 \(x\in P\) 都满足 \(x\leq x\)；
\item 反对称性：\(x\leq y\) 且 \(y\leq x\) 蕴含 \(x=y\)；
\item 传递性：\(x\leq y\) 且 \(y\leq z\) 蕴含 \(x\leq z\)。
\end{enumerate}

若 \(x\leq y\) 且 \(x\ne y\)，写作 \(x<y\)。子集 \(C\subseteq P\) 若任意两个元素 \(c,d\) 都满足 \(c\leq d\) 或 \(d\leq c\)，称为\term[lian]{链}[chain]。若整个 \(P\) 都是链，这个偏序又称\term[quanxu]{全序}[total order]。
\end{definition}

一个集合 \(X\) 的全部子集组成的集合称为幂集，记为 \(\mathcal P(X)\)。包含关系 \(\subseteq\) 是 \(\mathcal P(X)\) 上的偏序：每个子集包含于自身，相互包含的两个子集相等，包含关系还可以传递。但两个单点集 \(\{1\}\) 和 \(\{2\}\) 互不包含，所以 \(\mathcal P\bigl(\{1,2\}\bigr)\) 不是全序集；其中 \(\varnothing\subseteq\{1\}\subseteq\{1,2\}\) 则是一条链。

\begin{definition}\label{b1:def:A-upper-maximal}
设 \((P,\leq)\) 为偏序集，\(C\subseteq P\)。若 \(u\in P\) 满足对每个 \(c\in C\) 都有 \(c\leq u\)，则称 \(u\) 为 \(C\) 在 \(P\) 中的\term[shangjie]{上界}[upper bound]。若 \(m\in P\) 满足：对每个 \(x\in P\)，关系 \(m\leq x\) 都蕴含 \(x=m\)，则称 \(m\) 为\term[jidayuan]{极大元}[maximal element]；这表示在 \(P\) 中不能再严格增大 \(m\)。若 \(m\) 还满足对每个 \(x\in P\) 都有 \(x\leq m\)，则称 \(m\) 为\term[zuidayuan]{最大元}[greatest element]。
\end{definition}

最大元必须大于等于所有元素；极大元只要求没有比它更大的元素，仍可与其他元素不可比。因此最大元必为极大元，反之不成立。只由 \(\{1\}\)、\(\{2\}\) 组成的偏序集有两个极大元，却没有最大元；若最大元存在，则由反对称性知道它只能有一个。上界则是相对于指定子集而言的，也不必属于该子集。例如在 \(\mathcal P(\{1,2\})\) 中，\(\{1,2\}\) 是 \(\{\{1\},\{2\}\}\) 的上界，但不是这个两元素子集的成员。

\ProseParagraphBreak
例如有限初始段 \(F_n=\{0,\ldots,n\}\) 构成 \(\mathcal P(\mathbb N)\) 中的一条链，它们的并 \(\mathbb N\) 是上界，却不是其中任何一个 \(F_n\)。如果改在“\(\mathbb N\) 的全部有限子集”这个偏序集中讨论，同一条链就没有上界，因为任何包含全部 \(F_n\) 的集合都是无限的。上界属于哪一个偏序集，是不能省略的条件。

\begin{theorem}[Zorn]\label{b1:thm:A-zorn}
设 \((P,\leq)\) 为非空偏序集。若 \(P\) 中的每条链在 \(P\) 中都有上界，则 \(P\) 有极大元。
\end{theorem}

\begin{proofsketch}
证明的想法是：若没有极大元，就沿一个良序不断严格增大；但一个集合容不下任意长的这种序列。以下用序数表示良序的次序类型，并用超限递归从此前阶段的数据定义下一元素；这与证明已有性质成立的超限归纳不同，二者的次序语言将在后面说明。应用 Zorn 引理时只需检查定理所列的链上界条件，无须使用下面构造中的序数技巧。

这里调用 Hartogs 引理\footnote{哈托格斯（Friedrich Hartogs，1874-1943），德国数学家，研究集合论与多复变函数；以不能嵌入给定集合的序数构造著称。}与超限递归定理；两者在不使用选择公理的通常集合论中成立。序型和超限递归所需的集合论背景见\cite[Chapters 5--6]{Johnstone1987LogicSetTheory}；Hartogs 引理及其证明见\cite[Lemma 7.1]{Johnstone1987LogicSetTheory}。

由选择公理，为 \(P\) 的每个非空子集固定一个选取元素的方法。假设 \(P\) 没有极大元。给定一条链 \(C\)，定理的假设先给出它在 \(P\) 中的上界 \(u\)；因为 \(u\) 不是极大元，还存在 \(v>u\)。于是对每个 \(c\in C\)，都有 \(c\leq u<v\)，所以 \(v\) 严格大于这条链的全部成员。固定选择函数后，可以不断从这些元素中选择下一个，极限阶段则对此前的整条链作同样选择。

这种过程不能持续到一个足够长的序数。Hartogs 引理断言：对任意集合 \(P\)，存在不能单射到 \(P\) 的序数 \(\eta\)。其理由是，\(P\) 的子集上的全部良序关系组成一个集合；把每个这样的关系送到它的序型，再由替换公理收集这些序型，得到一个序数集合。取比其中每个序数都大的 \(\eta\)。若 \(\eta\) 可单射到 \(P\)，其像就会带有序型为 \(\eta\) 的良序，矛盾。按超限递归，对每个 \(\alpha<\eta\)，此前选出的 \(\{x_\beta\mid\beta<\alpha\}\) 是链；从严格大于这条链全部成员的非空集合中选取 \(x_\alpha\)。不同阶段的值严格增大，就得到 \(\eta\to P\) 的单射，矛盾。因此极大元存在。

这里的选择公理用于统一选取下一元素，Hartogs 引理保证过程必须停止；超限递归定理保证上述逐阶段的选取构成一个函数。下文只介绍理解这些步骤所需的序数语言。
\end{proofsketch}

这个结论称为\term[Zorn-yinli]{Zorn 引理}[Zorn's lemma]。它的假设并不是“任意两个对象都有共同上界”，也不是“每条有限链都有上界”，而是每条链都有上界，包括无限链。后面所说的有向集只要求任意两个元素有共同上界，这不能代替 Zorn 的假设：\(\mathbb N\) 的任意两个有限子集都有有限的并作共同上界，但上面的无限链 \((F_n)\) 在有限子集的偏序集中没有上界。结论只保证存在一个极大元，并不保证唯一，更不保证有最大元。偏序集 \(P\) 中的每个元素都是空链的上界，所以在检查条件时，先找出 \(P\) 的一个元素即可处理空链。

\ProseParagraphBreak
在通常的集合论公理体系中，Zorn 引理与选择公理等价。上面的概要说明选择公理如何推出 Zorn 引理；反向推导可以直接用部分选择函数完成，而且能说明为什么极大性与“同时选取”有关。

\begin{proposition}\label{b1:prop:A-zorn-implies-choice}
若接受 Zorn 引理，则每个由非空集合组成的集合族 \((X_i)_{i\in I}\) 都有选择函数。
\end{proposition}

\begin{proof}
考虑所有部分选择函数 \(c:J\rightarrow\bigcup_{i\in I}X_i\)，其中 \(J\subseteq I\)，且每个 \(i\in J\) 都有 \(c(i)\in X_i\)。把函数看成有序对的集合后，这些对象组成一个集合，记为 \(P\)。规定 \(c\leq d\) 当且仅当 \(d\) 延拓 \(c\)，也就是 \(\operatorname{dom}(c)\subseteq\operatorname{dom}(d)\) 且二者在 \(\operatorname{dom}(c)\) 上相等。这个关系是偏序；空函数属于 \(P\)，所以 \(P\) 非空。

设 \(\mathcal C\) 是 \(P\) 中的一条非空链。把链中函数的图像取并，得到 \(u=\bigcup_{c\in\mathcal C}c\)。若某个指标同时属于两个链成员的定义域，较大的函数延拓较小的函数，所以二者给出的值相同。因此 \(u\) 确实是函数，其定义域为各定义域的并。每个值都来自某个部分选择函数，故仍属于相应的 \(X_i\)。于是 \(u\in P\)，并且延拓每个链成员，是链的上界。空链用空函数作上界。

由\cref{b1:thm:A-zorn}，存在极大的部分选择函数 \(c:J\rightarrow\bigcup_iX_i\)。若 \(J\ne I\)，取一个 \(i_0\in I\setminus J\)。此时只面对这一个确定的非空集合 \(X_{i_0}\)；由它非空得到存在 \(x_0\in X_{i_0}\)，是普通存在量词的使用。把 \(c\) 扩充为在 \(i_0\) 处取值 \(x_0\) 的函数，便得到严格大于 \(c\) 的部分选择函数，矛盾。所以 \(J=I\)，即 \(c\) 是所求选择函数。最后一步没有预先对任意族作同时选择，因而没有循环使用待证的选择公理。
\end{proof}

\subsection{极大对象与基的扩充}
\label{b1:subsec:A01:02}

Zorn 引理本身只产生极大元。要得到某种代数对象，还须说明两件事：一条链的上界仍满足所要求的代数条件；一个不能再扩大的对象为什么已经达到目标。基的扩充把这两点表现得很清楚。检查链的并是否线性无关时，关键是一条线性关系只涉及有限个向量；这有限多个向量可以同时放进某一个链成员中，便能用该成员的线性无关性排除关系。

\begin{proposition}\label{b1:prop:A-basis-extension}
设 \(V\) 是域 \(K\) 上的向量空间。每个线性无关子集 \(A_0\subseteq V\) 都包含于 \(V\) 的一组基。特别地，每个向量空间都有基。
\end{proposition}

\begin{proof}
设 \(P\) 为所有包含 \(A_0\) 的线性无关子集组成的集合，按包含关系排序。这里若任何有限个互不相同的向量之间的线性关系都只能具有全零系数，就称这族向量线性无关。集合 \(A_0\) 属于 \(P\)，所以 \(P\) 非空。

设 \(\mathcal C\subseteq P\) 是非空链，令 \(U=\bigcup_{A\in\mathcal C}A\)。要检查 \(U\) 线性无关，取 \(U\) 中互不相同的 \(u_1,\ldots,u_r\)，并设
\[
a_1u_1+\cdots+a_ru_r=0,\qquad a_i\in K.
\]
每个 \(u_i\) 属于某个链成员 \(A_i\)。有限多个链成员之间两两可比，逐次比较可在 \(A_1,\ldots,A_r\) 中取出一个包含其余各集的成员 \(A\)。于是所有 \(u_i\) 都属于这个线性无关集，故 \(a_1=\cdots=a_r=0\)。这证明 \(U\) 线性无关。它又包含 \(A_0\)，所以 \(U\in P\)，并为链提供上界。空链则用 \(A_0\) 作上界。

由\cref{b1:thm:A-zorn}，取 \(P\) 中的极大元 \(B\)。剩下要证明 \(B\) 生成 \(V\)。若有 \(v\in V\) 不属于 \(B\) 的线性生成空间，则 \(B\cup\{v\}\) 仍线性无关。事实上，若
\[
av+b_1v_1+\cdots+b_sv_s=0,
\qquad v_1,\ldots,v_s\in B,
\]
其中 \(v_i\) 互不相同，则 \(a\ne0\) 会把 \(v\) 表成 \(B\) 的线性组合，与选择 \(v\) 的条件矛盾。因此 \(a=0\)，再由 \(B\) 线性无关得全部 \(b_i=0\)。于是 \(B\cup\{v\}\) 是 \(P\) 中严格包含 \(B\) 的成员，违反极大性。故 \(B\) 生成 \(V\)，与其线性无关性合起来便说明 \(B\) 是一组基。取 \(A_0=\varnothing\)，得到基的存在性。
\end{proof}

链的条件在这个证明中确有作用。设 \(e_1,e_2\) 为 \(K^2\) 的标准基，\(\{e_1\}\) 与 \(\{e_1+e_2,e_2\}\) 都线性无关。对第二个集合，关系 \(a(e_1+e_2)+be_2=0\) 给出 \(ae_1+(a+b)e_2=0\)，故 \(a=0\) 且 \(a+b=0\)，从而 \(b=0\)。但这两个集合的并因 \((e_1+e_2)-e_1-e_2=0\) 而相关。链使有限个向量能同时落在一个成员中；只有各个集合分别无关，并不足够。

\ProseParagraphBreak
这个证明也解释了为什么使用“极大”而不要求“最大”。例如在上一例的 \(K^2\) 中，不可能有一个线性无关集同时包含那两个无关集，因而不存在包含全部线性无关集的最大者。我们只需一组不能再添入向量的无关集，然后用上面的线性关系证明它必生成全空间。

\ProseParagraphBreak
上界检查中的有限性可以抽象成下面的引理。它同时适用于线性无关和代数独立，因为一条线性关系或一个多项式关系只涉及有限多个元素。

\begin{lemma}\label{b1:lem:finite-character-union}
设 \(X\) 为集合，\(\mathscr F\subseteq\mathcal P(X)\) 是一族子集，满足：对每个 \(A\subseteq X\)，\(A\in\mathscr F\) 当且仅当 \(A\) 的每个有限子集都属于 \(\mathscr F\)。若 \(\mathscr F\) 非空，则每个 \(A_0\in\mathscr F\) 都包含于 \(\mathscr F\) 中某个按包含关系极大的成员。
\end{lemma}

\begin{proof}
考虑 \(\mathscr F\) 中包含 \(A_0\) 的成员所成的偏序集。它包含 \(A_0\)。对其中非空链 \(\mathcal C\) 令 \(U=\bigcup_{A\in\mathcal C}A\)。若 \(F\subseteq U\) 有限且非空，为 \(F\) 的每个元素取一个含它的链成员；在有限多个成员中取最大的一个，便知 \(F\) 包含于某个 \(A\in\mathcal C\)。由 \(A\in\mathscr F\) 及假设，\(F\in\mathscr F\)。空集也是任一 \(A\in\mathcal C\) 的有限子集，故同样属于 \(\mathscr F\)。因此 \(U\) 的每个有限子集都属于 \(\mathscr F\)，再用假设得到 \(U\in\mathscr F\)。它包含 \(A_0\)，给出链的上界；空链用 \(A_0\) 作上界。\cref{b1:thm:A-zorn}遂给出所需的极大成员。
\end{proof}

引理中的条件称为有限特征性质。这个“当且仅当”的两个方向分别用于两次推理：从链成员 \(A\in\mathscr F\) 得到它的有限子集 \(F\in\mathscr F\)，再从并 \(U\) 的所有有限子集都属于 \(\mathscr F\) 得到 \(U\in\mathscr F\)。只保留后一个方向，尚不能完成前一个推理。

\ProseParagraphBreak
理想或部分映射的存在性证明虽然也使用链的并，却应各自检查其条件，不能直接把上面的性质套过去。例如“是一个理想”不具有这里的有限特征性质，因为理想的有限子集通常不是理想。对理想链的并 \(U\)，应直接取 \(x,y\in U\)，把它们放到同一个链成员中，得到 \(x-y\in U\)；若 \(r\) 为底环元素，\(rx\) 也在包含 \(x\) 的成员中。真理想之链的并仍为真，是因为单位元若进入并，已经进入某个成员。对部分映射，则须像\cref{b1:prop:A-zorn-implies-choice}的证明那样先检查不同定义域上的取值相容。

\ProseParagraphBreak
正文中的几个用法分别保留了这些核查：\cref{b1:prop:maximal-ideal-exists}构造包含给定真理想的极大理想；\cref{b1:thm:transcendence-basis}通过代数独立子集构造超越基，并用\cref{b1:lem:independence-adjoining}把极大性转成剩余元素的代数性；\cref{b1:thm:embedding-extension}将相容的部分域嵌入取并，再用最小多项式证明极大嵌入不能停在较小的定义域；\cref{b1:lem:finite-discrete-product-compact}则对真滤子使用 Zorn 引理。这些应用各有其数学背景，而上面基扩充的例子已足以说明极大性论证的基本做法。

\ProseParagraphBreak
以上论证给出存在性。如果还要求写出基向量、极大理想的生成元或域元素的有效表示，就必须另有计算条件。有限维基扩充可以在给定生成组中逐个试添向量，并用线性方程检查无关性；一般向量空间的基存在性没有提供这样的有限计算。

\subsection{良序与超限归纳}
\label{b1:subsec:A01:03}

普通归纳法按照 \(0,1,2,\ldots\) 的次序逐项证明结论。它之所以有效，可以用最小反例来说明：若反例存在，取最小的一个，归纳假设便在它之前全部成立。要把这个理由用于其他次序，所需条件就是每个非空子集都有最小元。

\begin{definition}\label{b1:def:A-well-order}
设 \((I,\leq)\) 为全序集。若对任意非空子集 \(J\subseteq I\)，都存在 \(j_0\in J\)，使每个 \(j\in J\) 都有 \(j_0\leq j\)，则称这个次序为\term[liangxu]{良序}[well-order]，并称 \(I\) 为良序集。
\end{definition}

\(\mathbb N\) 的通常次序是良序；\(\mathbb Z\) 的通常次序是全序，却不是良序，因为 \(\mathbb Z\) 自身没有最小元。一个集合能否良序，还取决于给它安排什么次序。例如按
\[
0,\ 1,\ -1,\ 2,\ -2,\ 3,\ -3,\ \ldots
\]
排列 \(\mathbb Z\)，就得到另一种次序。在这一次序中，任意非空子集都在上述序列中第一次出现，因而有最小元。

\begin{claim}[良序原理]\label{b1:claim:A-well-order}
每个集合都可赋予一个良序。
\end{claim}

\begin{proof}
给定集合 \(X\)，考虑 \(X\) 的子集上全部良序。不能只按底集合包含关系排序，因为较大的集合上安排的次序未必与较小集合相符；即使保留旧次序，在旧元素之间不断插入新元素，也未必使链的并仍为良序。因此按初始段延拓排序：较大的良序保留较小的次序，并把新增元素全部放在原集合之后。空集上的良序使这个偏序集非空。

对其中一条非空链，把底集合与次序关系分别取并。任意两个元素都同属某个链成员，其先后次序由该成员决定；不同成员给出的次序相符，所以得到全序。要说明它仍为良序，对并中任一非空子集 \(S\)，取 \(s\in S\) 及含 \(s\) 的链成员 \(Y\)。\(S\cap Y\) 有最小元 \(m\)。对其他链成员 \(Z\)，若 \(Z\) 是 \(Y\) 的初始段，其元素已经在 \(Y\) 中；若 \(Y\) 是 \(Z\) 的初始段，则 \(Z\setminus Y\) 的全部元素都在 \(Y\) 的全部元素之后，特别在 \(s\) 之后。因此 \(S\) 中不属于 \(Y\) 的元素都大于 \(s\geq m\)，故 \(m\) 是整个 \(S\) 的最小元。每个链成员又都是所得良序的初始段，所以并给出链的上界；空链用空集上的良序作上界。

Zorn 引理给出极大的良序子集 \(Y\)。若 \(Y\ne X\)，把 \(X\setminus Y\) 中任一元素添到最后，便得到严格延拓，矛盾。因此 \(Y=X\)，得到了 \(X\) 上的良序。
\end{proof}

良序原理也与选择公理等价。选择公理经 Zorn 引理推出上面的结论；反过来，若把 \(\bigcup_iX_i\) 良序，则对每个非空 \(X_i\) 取它在该次序中的最小元，就得到选择函数。这里得到的是某个良序的存在性，集合原有的通常次序可以不是良序；存在性论证也没有给出求取该良序的算法。

\begin{proposition}\label{b1:prop:well-order-induction}
设 \(I\) 为良序集，\(P(i)\) 是关于 \(i\in I\) 的性质。若对每个 \(i\in I\)，只要所有 \(j<i\) 都满足 \(P(j)\)，就能推出 \(P(i)\)，则每个 \(i\in I\) 都满足 \(P(i)\)。此外，良序集中不存在无限严格递降序列。
\end{proposition}

\begin{proof}
若有不满足性质的元素，这些反例组成的非空子集有最小元 \(i\)。由最小性，所有 \(j<i\) 都满足性质，假设却推出 \(i\) 也满足性质，矛盾。特别地，在最小元素处，“所有更早元素都满足”是空条件，所以命题的假设也已经包含归纳起点。

若有 \(i_0>i_1>i_2>\cdots\)，取值集合 \(\{i_n\mid n\in\mathbb N\}\) 有最小元，设为 \(i_m\)。但 \(i_{m+1}<i_m\) 仍属于这个集合，矛盾。
\end{proof}

这个论证称为良序归纳，也称\term[chaoxian-guina]{超限归纳}[transfinite induction]，它证明的是一个已经定义的性质在每个位置都成立。若用更早阶段的数据逐步定义新的对象，则是超限递归；前面的 Zorn 证明概要使用的正是后者。与自然数归纳相比，良序归纳允许某个位置之前已经有无限多个位置。例如先排所有 \(a_0,a_1,a_2,\ldots\)，再排所有 \(b_0,b_1,b_2,\ldots\)，规定每个 \(a_i\) 都小于每个 \(b_j\)，得到的仍是良序。但 \(b_0\) 之前有全部 \(a_i\)，却没有一个紧邻 \(b_0\) 的前项。因此不能只检查“前一个位置推出下一个位置”，还须处理没有紧邻前项的情形。

\ProseParagraphBreak
通常用\term[xushu]{序数}[ordinal]表示良序的次序类型。若两个良序集之间存在保持次序的双射，就说它们具有相同的次序类型；序数只记录这样的排列方式，不记录元素的原来名称。有限良序的次序类型为 \(0,1,2,\ldots\)，其中 \(0\) 对应空集；自然数通常次序的类型记为 \(\omega\)。在全部自然数之后再添一个新位置，得到 \(\omega+1\)；上面两串元素的次序类型则为 \(\omega+\omega\)。\(\omega\)、\(\omega+1\) 与 \(\omega+\omega\) 所对应的良序集都是可数集，却具有不同的次序类型。

\ProseParagraphBreak
序数 \(\alpha\) 的各个位置由 \(\beta<\alpha\) 编号。\(0\) 是起始位置；若某位置 \(\alpha\) 有紧邻前项 \(\beta\)，便写 \(\alpha=\beta+1\)，称为后继位置。非零而没有紧邻前项的位置称为极限位置，\(\omega\) 是第一个例子。在递增的代数构造中，起始步给出初始对象，后继步添加一个元素，极限步则通常取以前各层的并。要用良序归纳证明各阶段的对象具有某个性质，须验证起始步，以及任意后继步、极限步中该性质的保持。

\ProseParagraphBreak
例如\cref{b1:thm:free-abelian-subgroup}把一组基写成 \((e_\alpha)_{\alpha<\kappa}\)，令 \(F_\alpha\) 由 \(e_\beta\)（\(\beta<\alpha\)）生成。每个群元素只使用有限多个基向量，因此在极限位置之前出现的元素，其全部非零坐标早已同属某一个较早阶段。正因如此，若 \(\lambda\) 是极限位置，则
\[
F_\lambda=\bigcup_{\alpha<\lambda}F_\alpha.
\]
右边显然包含于左边。反过来，一个 \(x\in F_\lambda\) 只含有限多个基向量；这些非零系数所在的指标称为它的支集。若支集非空，取其中最大的指标 \(\beta<\lambda\)。因 \(\lambda\) 不是后继位置，有 \(\beta+1<\lambda\)，故 \(x\in F_{\beta+1}\) 已属于右边。零元也属于右边。于是极限步不会产生必须同时使用无限多个基向量的新元素；这个等式来自有限线性组合的定义，与分析中的级数收敛无关。

\ProseParagraphBreak
上面的 \(F_\alpha\) 已可直接由基和指标定义，无须递归求出；若此前的对象给出以后阶段的定义规则，才需要超限递归来保证这些阶段组成一个函数。一般超限递归定理以及序数的公理化构造属于集合论，相关背景见\cite[Chapters 5--6]{Johnstone1987LogicSetTheory}；本节只说明正文和下一小节用到的次序语言。

\subsection{无限集合与有限支集计数}
\label{b1:subsec:A01:05}

比较无限基时，每个基向量仍只涉及有限多个坐标，但要估计所有支集合起来的大小。本小节所需的两个计数结论是：一个无限集合的有限子集总数与它本身一样大；由无限多个有限集合组成的族，其并的大小不超过指标集的大小。这些结论用于任意秩自由交换群、无限维向量空间和无限超越基的基数比较。

\ProseParagraphBreak
有限集合可以数出元素个数，无限集合则用双射来比较大小：若 \(A\) 与 \(B\) 之间存在双射，称它们等势，写作 \(\lvert A\rvert=\lvert B\rvert\)；若存在单射 \(A\rightarrow B\)，写作 \(\lvert A\rvert\leq\lvert B\rvert\)。表示这种大小的数称为\term[jishu]{基数}[cardinal]。可数无限集与 \(\mathbb N\) 等势，它们的基数记为 \(\aleph_0\)。

\ProseParagraphBreak
例如 \(n\mapsto2n\) 是自然数集到偶数集的双射。无限集合的真子集可以与原集合等势，所以不能再以“删去一些元素”断言基数严格减小。下一结果说明，如果双方分别能单射到对方，那么它们确实等势。这个结论不需要选择公理。

\begin{proposition}[Cantor--Bernstein]\label{b1:prop:cantor-bernstein}
若集合 \(A,B\) 之间分别存在单射 \(f:A\rightarrow B\)、\(g:B\rightarrow A\)，则 \(A,B\) 之间存在双射。
\end{proposition}

\begin{proof}
把 \(A\) 与 \(B\) 中的元素想成两类点，\(f\) 的箭头从 \(A\) 指向 \(B\)，\(g\) 的箭头从 \(B\) 指回 \(A\)。单射 \(g\) 给出 \(g(B)\) 到 \(B\) 的逆映射，因此对属于 \(g(B)\) 的元素，可以沿 \(g\) 倒着走；不属于 \(g(B)\) 的元素却没有这种选择。令
\[
A_0=A\setminus g(B),\qquad
A_{n+1}=g\bigl(f(A_n)\bigr),\qquad
C=\bigcup_{n\geq0}A_n.
\]
集合 \(C\) 收集从这些没有 \(g\)-原像的元素出发，反复经过 \(f\) 再经过 \(g\) 所到达的元素。定义 \(h:A\rightarrow B\) 如下：
\[
h(a)=
\begin{cases}
f(a),&a\in C,\\
g^{-1}(a),&a\in A\setminus C.
\end{cases}
\]
第二行有意义，因为 \(A\setminus C\subseteq g(B)\)，且 \(g\) 单射。

由定义，\(g\bigl(f(C)\bigr)=\bigcup_{n\geq1}A_n=C\setminus A_0\)；这里 \(A_0\) 与所有 \(A_n\)（\(n\geq1\)）不交，因为后者都包含于 \(g(B)\)。所以 \(g\) 将 \(f(C)\) 双射到 \(C\setminus A_0\)。若 \(b\in B\setminus f(C)\)，则 \(g(b)\notin A_0\)；假如还有 \(g(b)\in C\)，便有 \(g(b)\in C\setminus A_0=g(f(C))\)。由 \(g\) 的单射性，这会推出 \(b\in f(C)\)，矛盾。因此 \(g(b)\in A\setminus C\)。反过来，若 \(a\in A\setminus C\)，它属于 \(g(B)\)，而 \(g^{-1}(a)\) 不能属于 \(f(C)\)，否则 \(a\in g(f(C))\subseteq C\)。这就说明 \(g^{-1}\) 将 \(A\setminus C\) 双射到 \(B\setminus f(C)\)。

于是 \(h\) 在 \(C\) 上双射到 \(f(C)\)，在 \(A\setminus C\) 上双射到 \(B\setminus f(C)\)。两部分的像不交且合为 \(B\)，故 \(h\) 是双射。
\end{proof}

以上结论称为\term[Cantor-Bernstein-dingli]{Cantor--Bernstein 定理}[Cantor--Bernstein theorem]。\footnote{康托尔（Georg Cantor，1845-1918），德国数学家，创立集合论，建立无穷集合的基数与序数理论。}\footnote{费利克斯·伯恩施坦（Felix Bernstein，1878-1956），德国数学家，证明两个集合互有单射时存在双射。}使用时不必真的写出所得双射；分别构造两个方向的单射，就足以证明两个集合大小相同。

\ProseParagraphBreak
先看可数情形。将自然数对按 \(\max(m,n)\) 分组，同组内先比较第一坐标，再比较第二坐标。行是第一坐标，列是第二坐标；前九对的编号为
\[
\begin{array}{c|ccc}
\alpha\backslash\beta & 0 & 1 & 2 \\ \hline
0 & 0 & 1 & 4 \\
1 & 2 & 3 & 5 \\
2 & 6 & 7 & 8
\end{array}
\]
这个次序先排完最大坐标为 \(0\) 的一组，再排最大坐标为 \(1\)、为 \(2\) 的各组，并非普通的逐行字典序。每组有限，每对都在某一组，于是可以把全部自然数对排成一列。这说明 \(\lvert\mathbb N\times\mathbb N\rvert=\aleph_0\)。反复使用这一计数，任意固定长度的自然数列也可数；把每个有限序列编码为“长度、该长度下的编号”这一自然数对，就得到全部有限序列的可数性。下面把这个现象推广到任意无限基数。

\ProseParagraphBreak
对一般无限集合，不能预先假定已经有自然数式的编号。证明将把集合良序，再比较每一对之前有多少个前驱，以此控制整个排列的大小。由良序原理，一个集合可以良序。在这个良序的次序类型以及比它更小的序数中，取与原集合等势的最小者，便得到表示其基数的标准序数。这样的序数称为初始序数。以下将无限基数 \(\kappa\) 与这个标准序数认同。因此 \(\alpha<\kappa\) 时，\(\alpha\) 的基数严格小于 \(\kappa\)；否则 \(\alpha\) 就是更小的等势序数。序数可以比较大小，而且任意非空的序数集合都有最小元，所以出现某个反例时，可以在不超过它的基数中选取最小反例。这些是本小节使用的序数基本性质。

\begin{proposition}\label{b1:prop:finite-support-cardinality}
采用选择公理。若 \(B\) 是无限集合，基数为 \(\kappa\)，则 \(B\times B\)、\(B\times\mathbb N\)、\(B\) 的全部有限序列及全部有限子集，都有基数 \(\kappa\)。由 \(\kappa\) 个有限集合组成的集合族，其并的基数至多为 \(\kappa\)。
\end{proposition}

\begin{proof}
先证明 \(\lvert\kappa\times\kappa\rvert=\kappa\)。假定某个无限基数不满足等式，取最小的反例 \(\kappa\)。把 \(\kappa\times\kappa\) 中的有序对先按 \(\max(\alpha,\beta)\) 的大小排列；最大值相同时，再先比较第一坐标、后比较第二坐标。所得次序是良序：在任意非空的有序对子集中，可以依次取最小的最大值、最小的第一坐标和最小的第二坐标，从而得到最小的一对。

固定其中一对 \((\alpha,\beta)\)，令 \(\gamma=\max(\alpha,\beta)\)。在它前面的所有有序对，都属于 \((\gamma+1)\times(\gamma+1)\)。记 \(\mu=\lvert\gamma+1\rvert\)，则 \(\mu<\kappa\)。这里还需说明添上最后一个位置不会达到 \(\kappa\)：若 \(\gamma\) 有限则显然；若 \(\gamma\) 无限，它包含自然数的全部位置，将新添元素送到 \(0\)，将每个自然数 \(n\) 送到 \(n+1\)，将其余元素固定，就给出 \(\gamma+1\) 到 \(\gamma\) 的双射。因此 \(\lvert\gamma+1\rvert=\lvert\gamma\rvert<\kappa\)。这里相等的是基数；作为序数，\(\gamma+1\) 仍严格大于 \(\gamma\)，上述双射并不保持原有次序。

若 \(\mu\) 有限，\((\gamma+1)\times(\gamma+1)\) 仍有限，其基数小于 \(\kappa\)；若 \(\mu\) 无限，由反例 \(\kappa\) 的最小性，这个积集的基数仍为 \(\mu<\kappa\)。所以在上述良序中，每一对之前的初始段都具有小于 \(\kappa\) 的基数。设这个良序的次序类型为 \(\theta\)。若 \(\theta>\kappa\)，其中便存在序号为 \(\kappa\) 的位置，而它之前恰有与 \(\kappa\) 等势的全部位置，矛盾。故 \(\theta\leq\kappa\)，把各对送到它们的序号，就得到 \(\kappa\times\kappa\) 到 \(\kappa\) 的单射。反向单射为 \(\alpha\mapsto(\alpha,0)\)，由\cref{b1:prop:cantor-bernstein}得到等势，排除了反例。

将 \(B\) 与 \(\kappa\) 通过双射认同，就得到 \(\lvert B\times B\rvert=\kappa\)。因为无限初始序数 \(\kappa\) 包含所有自然数，\(\mathbb N\) 能单射到 \(B\)，故
\[
\kappa\leq\lvert B\times\mathbb N\rvert
\leq\lvert B\times B\rvert=\kappa.
\]
再用\cref{b1:prop:cantor-bernstein}，得到 \(\lvert B\times\mathbb N\rvert=\kappa\)。对长度 \(n\geq1\) 作有限归纳，得到 \(\lvert B^n\rvert=\kappa\)。对每个长度 \(n\)，单射 \(B^n\rightarrow B\) 所成的集合非空；用选择公理同时从这些集合中各选一个单射。这些编码只需存在，并不要求有典范选法。于是可把全部有限序列映入 \(\mathbb N\times B\)：第一个坐标记录长度，第二个坐标记录该长度下的编码；空序列也可选定一个编码。长度一的序列给出反向下界，所以全部有限序列的基数为 \(\kappa\)。

在 \(B\) 上固定良序，每个有限子集可唯一地按递增次序写成一个有限序列，因而有限子集的总数至多为 \(\kappa\)；单点子集又给出下界 \(\kappa\)，故两者相等。

最后，设 \((E_i)_{i\in I}\) 是一族有限集合，\(\lvert I\rvert=\kappa\)。为每个 \(E_i\) 选取一个单射到 \(\mathbb N\) 的编号，则带指标的集合 \(\bigl\{\,(i,x)\mid i\in I,\ x\in E_i\,\bigr\}\) 单射到 \(I\times\mathbb N\)，故基数至多为 \(\kappa\)。不带指标的并 \(\bigcup_iE_i\) 是这个集合的满射像；把带指标的集合良序，为并中的每个元素取最小的原像，就得到反向单射，证明并的基数至多为 \(\kappa\)。同时选定编号和良序的步骤都使用已经采用的选择公理。
\end{proof}

例如一族有限支集 \((S_c)_{c\in C}\) 若覆盖无限集合 \(B\)，那么上面的命题给出 \(\lvert B\rvert\leq\lvert C\rvert\)，这里 \(C\) 也必无限。这比“每个支集都有限”多了一步总量比较：有限支集的大小不必有共同上界，但其并仍不会超过无限指标集的基数。

\ProseParagraphBreak
例如设 \(B,C\) 为同一个 \(K\)-向量空间的两组基，把每个 \(c\in C\) 写成
\[
c=\sum_{b\in S_c}a_{b,c}b,
\qquad S_c\subseteq B\;\text{有限},\quad a_{b,c}\in K.
\]
若 \(b_0\in B\) 不属于 \(\bigcup_{c\in C}S_c\)，那么所有 \(c\in C\) 都属于 \(\operatorname{span}_K(B\setminus\{b_0\})\)。由于 \(B\) 线性无关，这个子空间不含 \(b_0\)，所以 \(C\) 不可能生成整个空间，矛盾。因此支集的并确实覆盖 \(B\)。当 \(B\) 无限时，\(C\) 也必无限，命题便给出 \(\lvert B\rvert\leq\lvert C\rvert\)；交换两组基的角色，再用 Cantor--Bernstein 定理得到等号。

\ProseParagraphBreak
这正是\cref{b1:prop:free-rank-invariant}中的无限计数；\cref{b1:prop:free-abelian-rank}对自由交换群使用整数坐标时，道理相同。\cref{b1:thm:transcendence-cardinality}还要先用有限超越次数估计控制每个有限支集所对应的元素数目，随后才作同样的无限计数。有限情形的交换论证与这里的基数比较各有职责，不能仅把“有限”一词删去便得到无限情形。
